\chapter{Training for Reasoning}
\minitoc

\begin{chapteroverview}
  \begin{itemize}
    \item Analyze Chain-of-Thought (CoT) and its 2025--2026 training methods.
    \item Master Process Reward Models (PRM) for step-level supervision.
    \item Implement test-time compute scaling and search (MCTS, Best-of-N).
    \item Understand Thinking Tokens (Quiet-STaR) for emergent reasoning.
  \end{itemize}
\end{chapteroverview}

\section{The Rise of Reasoning Models}

Models that break complex problems into explicit intermediate steps---the most significant capability development of 2025--2026. Reasoning models sacrifice token efficiency for accuracy on hard tasks.

\section{Chain-of-Thought (CoT)}

Prompting or training models to generate intermediate reasoning steps before the final answer. Few-shot CoT \parencite{wei2022chain} demonstrated step-by-step prompting dramatically improves multi-step arithmetic, symbolic reasoning, and commonsense tasks.

\begin{itemize}
  \item \textbf{Zero-shot CoT:} Append ``Let's think step by step'' to the prompt.
  \item \textbf{Few-shot CoT:} Provide example (question, reasoning chain, answer) triples in the prompt.
  \item \textbf{Self-consistency:} Sample $k$ CoT paths and take the majority vote answer. +5--15\% on GSM8K vs.\ single-sample CoT.
\end{itemize}

\section{Process Reward Models (PRM)}

Instead of rewarding only the final answer, PRMs assign a score to each reasoning step:
\begin{itemize}
  \item Denser supervision signal: $n$ steps $\rightarrow$ $n$ reward signals per example.
  \item Enables early pruning of incorrect branches during tree search.
  \item Used in OpenAI's o-series for competition math and coding \parencite{lightman2023lets}.
\end{itemize}

Annotation challenge: step-level labels are expensive. Solutions: automated labeling via outcome consistency (steps on paths leading to correct answers are labeled positive) and synthetic step decomposition from a stronger model.

\section{Test-Time Search}

Allocate more inference compute per query to find better solutions:
\begin{itemize}
  \item \textbf{Best-of-N:} Sample $N$ complete solutions; select via outcome RM or verifier.
  \item \textbf{Beam Search over CoT:} Expand $k$ partial chains at each step; prune with PRM.
  \item \textbf{Monte Carlo Tree Search (MCTS):} Full tree exploration with PRM scoring. Used in AlphaCode 2 and OpenAI o3.
\end{itemize}

\begin{keyinsight}[Test-Time Compute Scaling]
\textcite{snell2024scaling} showed that test-time compute can match training-time compute scaling: a model with 16$\times$ more test-time compute can match a model 16$\times$ larger. This creates a new performance axis independent of parameter count.
\end{keyinsight}

\section{TLT -- Accelerating Training}

MIT: a smaller ``drafter'' model predicts outputs; a larger verifier model confirms them \parencite{leviathan2023fast}. 70--210\% acceleration with preserved accuracy. The drafter is reusable for speculative decoding at inference.

\section{R1-Style Pure RL Training}

DeepSeek demonstrated that pure RL on verifiable tasks produces emergent reasoning without any SFT \parencite{guo2025deepseekr1}:
\begin{enumerate}
  \item Start from the base model (no instruction tuning).
  \item Apply GRPO (\hyperref[form:grpo]{formulation \S\ref*{form:grpo}}) with binary correctness reward on math/code tasks.
  \item Emergent behaviors: self-verification (``wait, let me reconsider''), extended thinking, strategy switching.
\end{enumerate}

\begin{warningbox}[Pure RL Training Instability]
R1-Zero (pure RL, no SFT cold start) exhibits language mixing, repetitive patterns, and poor readability. The DeepSeek production pipeline uses a cold-start SFT phase (a few thousand high-quality CoT examples) before GRPO to stabilize training.
\end{warningbox}
\section{Thinking Tokens (Quiet-STaR)}

Rather than forcing reasoning into a separate pre-response phase, Quiet-STaR \parencite{schulman2024quietstar} trains the model to generate ``internal monologue'' tokens at every position in the sequence:
\begin{itemize}
  \item \textbf{Divergent Thinking:} The model generates multiple parallel reasoning traces (thoughts) in the background.
  \item \textbf{Rationalization:} It selects the thought that most improves the prediction of the subsequent actual text tokens.
  \item \textbf{Token-Level RL:} Thoughts are rewarded based on how much they reduce the cross-entropy loss of the next real token.
\end{itemize}
This allows models to benefit from reasoning even on tasks where explicit Chain-of-Thought is not requested, and it creates a unified architecture for both ``fast'' (direct) and ``slow'' (reasoned) processing.

\begin{keyinsight}[The Thinking Token Trade-off]
Thinking tokens increase inference compute (latency and FLOPs) but can dramatically reduce the model size required to achieve a certain level of performance. A 7B model with 16 thinking tokens can match a 70B model on several reasoning-heavy benchmarks.
\end{keyinsight}
